\subsection{4. Đối tượng và phạm vi nghiên cứu}

\indent\indent Đối tượng nghiên cứu của đề tài là các mô hình học sâu đa phương thức dựa trên kiến trúc đồ thị, bao gồm các cơ chế tổng hợp thông tin giữa hình ảnh, văn bản và âm thanh. Cốt lõi của đề tài là tìm hiểu, xây dựng và đánh giá các mô hình GNN (đặc biệt là GAT) kết hợp với các kiến trúc mới như KAN (FastKAN), cùng với MLP và các mô hình máy học truyền thống trong vai trò bộ phân lớp. Đề tài hướng đến việc khảo sát cách các phương thức dữ liệu khác nhau có thể được biểu diễn và liên kết trong một cấu trúc đồ thị thống nhất để phục vụ nhiệm vụ phân loại văn hoá phi vật thể.\vspace{0.3cm}

Phạm vi nghiên cứu được giới hạn trong phạm vi dữ liệu video văn hoá được thu thập từ YouTube, tập trung vào các chủ đề như lễ hội dân gian và nghề thủ công truyền thống tại Việt Nam. Từ các video này, nghiên cứu chỉ khai thác ba loại dữ liệu chính: (1) khung hình trong video, (2) văn bản phụ đề hoặc mô tả, và (3) đặc trưng âm thanh. Các phân tích, thử nghiệm và đánh giá đều được thực hiện trong khuôn khổ bộ dữ liệu này, không mở rộng sang các loại hình di sản khác hoặc các nguồn dữ liệu ngoài phạm vi đã nêu.